\documentclass{article}
\usepackage{anyfontsize}
%\usepackage[UTF8]{ctex} % 如果需要支持中文
\usepackage{xeCJK} % XeLaTeX 推荐 xeCJK，LuaLaTeX 也可用

\usepackage{fontspec} 

% \setCJKmainfont{SourceHanSansCN-Light}[
%     BoldFont=SourceHanSansCN-Medium
% ]


\usepackage[a4paper, left=0.1in, right=0.1in, top=0.05in, bottom=0.05in]{geometry} % 设置四边页边距为0.1英寸{geometry} % 设置页边距为0.5英寸
\usepackage{enumitem} % 用于自定义列表
\usepackage{amsmath}  % 数学公式支持

\setlength{\parindent}{0pt} % 不缩进
\usepackage{etoolbox} % 用于列表操作
%调整类代码
\usepackage{comment}
\usepackage{tocloft} % 用于定制目录 样式
\usepackage{hyperref} %不需要目录盒子注释这
\usepackage{ifthen}
\usepackage{tikz}
\usepackage{titletoc}
\usepackage{graphicx}
\usepackage{amssymb}  % 提供数学符号，包括\therefore
%目录设定
%快捷键 CMD + / 注释
%  \renewcommand{\cftdot}{} % 隐藏目录中的页码
%  \cftpagenumbersoff{section}
%  \cftpagenumbersoff{subsection}
%  \makeatletter
%  \renewcommand{\@pnumwidth}{0em}  % 设置页码宽度为0
%  \renewcommand{\@tocrmarg}{0em}   % 设置右边距为0
%  \makeatother
 


\usepackage{environ}

\newif\ifshowcontent  % 定义一个条件判断变量
\showcontenttrue    % 设置为 false，隐藏所有 question 环境的内容

\newcounter{question} % 题目计数器
\setcounter{question}{0}
\NewEnviron{question}{%
    \ifshowcontent
        \par\stepcounter{question}\textbf{\thequestion.} \BODY\par % 如果显示内容，则输出
    \fi
}

\NewEnviron{choices}{%
    \ifshowcontent
        \begin{enumerate}[label=\Alph*., leftmargin=*]
            \BODY
        \end{enumerate}\par
    \fi
}

\NewEnviron{correctanswer}{%
    \ifshowcontent
        \textbf{答案:}\BODY\par
    \fi
}



\newenvironment{explanation}
{\textbf{解析:}\par}
{\par}



% 新增考察方面命令
\newcommand{\checklist}[1]{
    \textbf{考察:} #1 \par % 在题目下方显示考察方面
    \addcontentsline{toc}{subsection}{题号\thequestion 考察: #1} % 将考察内容添加到目录
    \vspace{2em} % 添加1em的垂直空间
}

\graphicspath{{images/}} %统一


\begin{document}
 %\fontsize{22}{31}\selectfont
 \tableofcontents % 添加目录
\newpage
%\pagestyle{empty}




\section{考点1:回归对冲和主成分分析}
\setcounter{question}{0}



\begin{question}
    Assume that a trader is making a relative value trade, selling a U.S. Treasury bond and
    correspondingly purchasing a U.S. Treasury TIPS. Based on the current spread between the two
    securities, the trader shorts \$120 million of the nominal bond and purchases \$105.6 million of
    TIPS. The trader then starts to question the amount of the hedge due to changes in yields on
    TIPS in relation to nominal bonds. He runs a regression and determines from the output that the
    nominal yield changes by 1.0352 basis points per basis point change in the real yield. Would
    the trader adjust the hedge, and if so, by how much?
    \\假设一个交易员正在做一个相对价值交易，卖出一张美国国债，同时购买一张美国国债TIPS。根据当前两个证券之间的利差，交易员卖出\$120M名义债券，买入\$105.6M TIPS。然后，交易员开始质疑TIPS收益率变化与名义债券收益率变化之间的关系。他进行回归分析，并从输出中确定名义收益率每变化1个基点，实际收益率变化1.0352个基点。交易员是否会调整对冲，如果是，调整多少？
\end{question}

\begin{choices}
    \item No.
    \item Yes, by \$3.12 million (purchase additional TIPS)\\
        是的，需要额外购买\$3.12M的TIPS。
    \item Yes, by \$3.15 million (sell a portion of the TIPS)\\
        是的，需要卖出部分TIPS。
    \item Yes, by \$2.88 million (Purchase additional TIPS)\\
        是的，需要额外购买\$2.88M的TIPS。
\end{choices}

\begin{correctanswer}
    B
\end{correctanswer}

\begin{explanation}
  
    \textbf{步骤1：理解对冲条件}
        \begin{itemize}
            \item 完全对冲需要：名义债券的美元久期变化 = TIPS的美元久期变化
            \item 即：$MD_n \times P_n \times \Delta y_n = MD_r \times P_r \times \Delta y_r$
        \end{itemize}

    \textbf{步骤2：原始对冲假设}
        \begin{itemize}
            \item 原始头寸：卖空\$120M名义债券，买入\$105.6M TIPS
            \item 原始假设：$\Delta y_n = \Delta y_r$（名义收益率变化等于实际收益率变化）
            \item 因此：$MD_n \times 120 = MD_r \times 105.6$
        \end{itemize}

    \textbf{步骤3：根据回归结果调整}
        \begin{itemize}
            \item 回归显示：$\Delta y_n = 1.0352 \times \Delta y_r$
            \item 新的对冲条件：
            \[
            MD_n \times 120 \times \Delta y_n = MD_r \times P_r \times \Delta y_r
            \]
            \item 代入$\Delta y_n = 1.0352 \times \Delta y_r$：
            \[
            MD_n \times 120 \times 1.0352 \times \Delta y_r = MD_r \times P_r \times \Delta y_r
            \]
            \item 化简：$MD_n \times 120 \times 1.0352 = MD_r \times P_r$
        \end{itemize}

    \textbf{步骤4：计算新的TIPS头寸}
        \begin{itemize}
            \item 从步骤2：$MD_n = \frac{105.6}{120} \times MD_r$
            \item 代入步骤3的结果：
            \[
            \frac{105.6}{120} \times MD_r \times 120 \times 1.0352 = MD_r \times P_r
            \]
            \item 化简：$105.6 \times 1.0352 = P_r$
            \item 因此：$P_r = 108.72$ million
        \end{itemize}

    \textbf{步骤5：计算调整金额}
        \begin{itemize}
            \item 需要增加的TIPS头寸：$108.72 - 105.6 = 3.12$ million
            \item 因此需要额外购买\$3.12M的TIPS。
        \end{itemize}

\end{explanation}
\checklist{
    掌握回归对冲（根据收益率变化关系调整对冲比例）
}


\begin{question}
    If a trader is creating a fixed income hedge, which hedging methodology would be least
    effective if the trader is concerned about the dispersion of the change in the corporate bond yield
    for a particular change in the government bond yield?
    \\假设一个交易员正在创建一个固定收益对冲，如果交易员担心公司债券收益率变化与政府债券收益率变化之间的分散性，哪种对冲方法最无效？
\end{question}

\begin{choices}
    \item One-variable regression hedge.\\
        单变量回归对冲。
    \item DV01 hedge\\
        DV01对冲。
    \item Two-variable regression hedge.\\
        双变量回归对冲。
    \item Principal components hedge.\\
        主成分对冲。
\end{choices}

\begin{correctanswer}
    B
\end{correctanswer}

\begin{explanation}
    \textbf{A选项错误。}
    \begin{itemize}
        \item 单变量回归对冲可以捕捉公司债和政府债收益率变化之间的关系。
    \end{itemize}
    \textbf{B选项正确。}
    \begin{itemize}
        \item DV01对冲假设公司债和政府债的收益率变化幅度相同，无法处理两者之间的分散性。具体来说：
        \begin{itemize}
            \item DV01对冲基于收益率平行移动的假设
            \item 无法反映信用利差的变化
            \item 忽略了收益率曲线的非平行移动
        \end{itemize}
    \end{itemize}
    \textbf{C选项错误。}
    \begin{itemize}
        \item 双变量回归对冲可以更好地捕捉收益率变化的分散性。
    \end{itemize}
    \textbf{D选项错误。}
    \begin{itemize}
        \item 主成分对冲可以处理收益率曲线的各种变化形态。
    \end{itemize}
\end{explanation}
\checklist{
    理解对冲方法（DV01对冲无法处理收益率分散性）
}




\newpage


\end{document}